\documentclass[10pt,a4paper]{amsart}
\usepackage[margin=19mm]{geometry}
\usepackage{amsmath,amssymb,mathtools}
\usepackage[scr=rsfs,cal=euler]{mathalfa}
\usepackage{xfrac}
\usepackage[unicode,hidelinks,bookmarks=false]{hyperref}
\newcommand{\ie}{i.e.\ }
\newcommand{\cf}{cf.\ }
\newcommand{\resp}{respectively\ }
\newcommand{\Subsection}{\subsection}
\providecommand{\nobreakdash}{\nobreak\hbox{-}}
\setlength{\emergencystretch}{2em}
\multlinegap0pt
\usepackage{xcolor}
\usepackage[normalem]{ulem}
\usepackage{tikz}
\usepackage{amssymb}
\newtheorem{theorem}{Theorem}
\newtheorem{lemma}{Lemma}
\newcommand{\MBPD}{\textnormal{MBPD}}
\title{On modular balanced partition designs}
\author{S. Karthik, Peter J. Cameron, Krishnan Paramasivam}
\date{Source: arXiv:2607.18087v1 (CC BY 4.0)}
\begin{document}
\maketitle

\noindent\emph{Source and licence.} On modular balanced partition designs. Version arXiv:2607.18087v1 (CC BY 4.0), distributed by arXiv under the Creative Commons Attribution 4.0 International licence, http://creativecommons.org/licenses/by/4.0/, as recorded in the arXiv licence metadata for this exact version and shown on the abstract page https://arxiv.org/abs/2607.18087v1. Full licence notice: \texttt{licenses/arxiv-2607.18087-CC-BY-4.0.txt}.\par\smallskip

\noindent\textit{Editorial scope.} Counted unit: the article's lemma karlem (source0.tex lines 989--991). The article's complete original proof of it is delivered with it (source0.tex lines 992--1081, one \texttt{proof} environment of 90 lines). No mathematics is written, repaired or re-derived: the statement and the proof are the article's own lines, in the article's own notation, and no retained line is substituted or re-flowed.  Retained uncounted as the source-local context the proof needs: lines 607--609.  Nothing was omitted from any retained span, so the package carries no omission record.  No retained line contains a citation, so the wrapper carries no bibliography and no bibliography entry.  No retained line contains a figure environment, an image-inclusion command or drawing code, so no image is delivered and the package declares no asset dependency.  Editorial formatting only, in addition: an AMS \texttt{amsart} preamble that restates the article's own declarations and macros used by the retained lines, namely the environments \texttt{theorem}, \texttt{lemma} on separate counters, together with the standard packages those lines need.  The retained environments print in this document as Theorem 1, Lemma 1 in order of appearance, whereas the article prints the same items with the numbering of its own full document.  The complete native source of the article as distributed in the arXiv e-print for this exact version is bundled unchanged as \texttt{sources/205624/source0.tex}.
\begin{theorem}{\label{thmkar}}
      If $n\equiv  0 \bmod{4}$, then there exists a $(n^2,n ,n)$-$\MBPD$~over $\mathbb{Z}_{n^2}$.
\end{theorem}

\begin{lemma}\label{karlem}
   If $m\equiv 0\bmod{4}$ and $n=km,$ where $k$ is odd, then there exists an $m\times n$ subgroup magic rectangle over $\mathbb{Z}_{mn}$.
\end{lemma}

\begin{proof}
    The construction is given in three steps. 
    \\
    \noindent Step 1. Without loss of generality, assume $m \leq n$. By Theorem~\ref{thmkar},  let $\mathbb{M}$ be the $m \times m$ subgroup magic rectangle.\\
    \noindent Step 2. Let $\mathbb{A}$ be an $m \times n$ matrix with $n = km$, and $k$ is odd. 
Then $\mathbb{A}$ can be expressed as a concatenation of $k$ block matrices, each of size $m \times m$; that is,
\[
\mathbb{A} = \begin{bmatrix} 
\mathbb{A}_1 ~~& \mathbb{A}_2 ~~& \cdots & ~~\mathbb{A}_k
\end{bmatrix}
\]
Now, for each $i\in [1,k]$, we define the $m\times m$ matrix $\mathbb{A}_i$ such that \[\mathbb{A}_i=\begin{cases}
    \mathbb{M} & \text{if} ~i=1, \\
    m^2\mathbb{I}+ \mathbb{A}_{i-1} & \text{if}~ i \in [2,k].\end{cases},\]
where $\mathbb{I}$ is the $m\times m$ identity matrix.\\
\noindent Step 3. In this step, we prove that $\mathbb{A}$ forms a subgroup magic rectangle. Since $\mathbb{M}$ is a subgroup magic rectangle over $\mathbb{Z}_{m^2}$, the row-subgroup of $\mathbb{M}$ over $\mathbb{Z}_{m^2}$ is $\langle m\rangle$. We consider the expanded matrix $\mathbb{A}$ over $\mathbb{Z}_{k m^2}$, where $k$ is an odd integer and $m \equiv 0 \pmod{4}$, obtained by replicating and shifting $\mathbb{M}$ in blocks.

\noindent We claim that the following subset $H_1$ of the group $\mathbb{Z}_{k m^2}$,
\[
H_1 = \frac{(k-1
)k}{2}m^3 \bmod{km^2}+  \langle m \rangle   
\]
forms a subgroup under addition modulo $k m^2$.
Thus, \begin{align*}
    H_1 & = \{ikm\bmod{m^2}+ \frac{(k-1)k}{2}m^3\bmod{km^2}: i \in [1, m]\}\\
    & = \{ikm\bmod{m^2}:  i \in [1, m]\} = \{ikm\bmod{km^2}:  i \in [1, m]\}.
\end{align*}
\noindent Therefore, $H_1$ is  subgroup of order $m$. Similarly, the column sums of $\mathbb{A}$, are in the set $H_2$. Then, \begin{align*}
    H_2&= \{ilm\bmod{m^2}+(i-1)m^3\bmod{km^2} : i \in[1,k], l \in[1,m]\}
\end{align*} 
We first show that every element of \(\mathcal H_2\) belongs to the subgroup
\(\langle m\rangle\) of \(\mathbb Z_{km^2}\). Let
\[
h=ilm \bmod{m^2}+(i-1)m^3 \bmod{km^2}
\in \mathcal H_2.
\]
Since
\(
ilm \pmod{m^2}=m(il \pmod m),
\)
there exists an integer \(r\in[0,m-1]\) such that
\(
ilm \pmod{m^2}=rm.
\)
Hence
\(
h=rm+(i-1)m^3=m\bigl(r+(i-1)m^2\bigr).
\)
Therefore \(h\) is a multiple of \(m\), and thus
\(
h\in \langle m\rangle.
\)
Since \(h\) was arbitrary,
\(
\mathcal H_2\subseteq \langle m\rangle\). Now consider the subgroup generated by $\langle m\rangle$ in $\mathbb{Z}_{km^2}$. Since,
\[ o(m)=\frac{km^2}{\gcd(km^2,m)}=
\frac{km^2}{m}=
km.
\]
Thus \(\langle m\rangle\) is a subgroup of \(\mathbb Z_{km^2}\) of order \(km\). Observe also that
\(
m^3=m^2\cdot m,
\)
and hence
\(
m^3\in\langle m\rangle.
\)
Consequently,
\(
\langle m,m^3\rangle=\langle m\rangle.
\) Since the column sums of \(\mathbb A\) form a subgroup of order \(km\), we have
\(
|\mathcal H_2|=km.
\)
Combining this with
\(
\mathcal H_2\subseteq \langle m\rangle
~\text{and}~
|\langle m\rangle|=km,
\)
it follows that
\(
\mathcal H_2=\langle m\rangle.
\)
Therefore, 
\[H_2=\langle m\rangle=
\langle m,m^3\rangle=
\{tm \bmod{km^2}: t\in\mathbb Z\},
\]
\noindent which is a subgroup of \(\mathbb Z_{km^2}\) of order \(km\). This completes the proof.\end{proof}

\end{document}
